\documentclass[11pt]{article}
\usepackage[T1]{fontenc}
\usepackage{lmodern,amsmath,amssymb,amsthm,mathtools}
\usepackage[margin=1in]{geometry}
\usepackage{microtype,enumitem,booktabs,xcolor}
\usepackage[numbers,sort&compress]{natbib}
\usepackage{hyperref}
\hypersetup{colorlinks=true,linkcolor=blue!45!black,citecolor=blue!45!black,
urlcolor=blue!45!black,pdftitle={Bernoulli extremizers for the independent Berry-Esseen ratio},pdfauthor={Logan Bell}}
\setlength{\emergencystretch}{2em}
\numberwithin{equation}{section}
\newtheorem{theorem}{Theorem}[section]
\newtheorem{proposition}[theorem]{Proposition}
\newtheorem{lemma}[theorem]{Lemma}
\newtheorem{corollary}[theorem]{Corollary}
\theoremstyle{remark}
\newtheorem{remark}[theorem]{Remark}
\DeclareMathOperator{\Var}{Var}
\DeclareMathOperator{\supp}{supp}
\DeclareMathOperator{\Bern}{Bern}
\DeclareMathOperator{\Bin}{Bin}
\DeclareMathOperator{\osc}{osc}
\DeclareMathOperator{\sinc}{sinc}
\newcommand{\E}{\mathbb E}
\newcommand{\Pp}{\mathbb P}
\newcommand{\R}{\mathbb R}
\newcommand{\Z}{\mathbb Z}
\newcommand{\ind}{\mathbf1}
\newcommand{\ce}{C_{\mathrm E}}
\newcommand{\Cind}{C_{\mathrm{ind}}}
\newcommand{\pe}{p_{\mathrm E}}
\newcommand{\phz}{\phi(0)}
\newcommand{\eps}{\varepsilon}
\newcommand{\ii}{\mathrm i}
\title{Bernoulli extremizers for the independent\\Berry--Esseen ratio}
\author{Logan Bell}
\date{6 October 2026}
\begin{document}
\maketitle
\begin{abstract}
We study the extremal structure of the Berry--Esseen ratio for independent
centered two-point (Bernoulli) random variables, allowing deterministic zero
summands. First, for every $n\ge1$, we prove that every array maximizing the
Berry--Esseen ratio among $n$-summand two-point arrays with a common support
diameter consists of $m\le n$ identically distributed nondegenerate two-point
summands and $n-m$ zero summands. Second, for at most two two-point summands
with arbitrary support diameters, we prove that the maximum ratio equals the
one-summand constant $C_1\approx0.3703515235$, attained only when one summand
is zero. Third, we determine the asymptotic behavior of the optimal active
count $m_N$ in the binomial envelope $P_N=\max_{1\le m\le N}B_m$: a uniform
seventh-order Edgeworth--Euler--Maclaurin expansion and a quadratic-field
separation argument in $\mathbb Q(\sqrt{10})$ prove that for all sufficiently
large $N$, $m_N$ is unique, the number of zero summands $N-m_N$ belongs to
$\{0,1,2,3,4\}$ with explicit limiting frequencies, and successive record
active sizes have gaps in $\{2,3,5\}$. Fourth, for two-point arrays with
large count variance $A=\sum_{i=1}^np_i(1-p_i)$ and small diameter deviations
$e_i=g_i/h_0-1$, we prove that replacing all diameters by their
variance-weighted mean $h_0$ strictly increases the ratio by at least
$\kappa\sigma^2/(\sigma+\eta)$, where $\sigma^2=\sum_{i=1}^np_i(1-p_i)e_i^2$
and $\eta=\max_i|e_i|$, and we construct an explicit family showing that the
rate $\sigma^2/(\sigma+\eta)$ is optimal in the absence of a lower bound on
the individual probabilities.
\end{abstract}
\tableofcontents

\section{Introduction and main results}

\subsection{Formulation of the extremal problem}

Let $X_1,\ldots,X_n$ be independent centered real random variables with
positive total variance and finite third absolute moments, and set
\begin{gather*}
 V=\sum_{i=1}^n\E X_i^2>0,\qquad R=\sum_{i=1}^n\E|X_i|^3<\infty,\\
 \mathcal C(X_1,\ldots,X_n)=\frac{V^{3/2}}R
   \sup_{y\in\R}\left|\Pp\left\{\sum_{i=1}^nX_i\le y\right\}-\Phi\!\left(\frac y{\sqrt V}\right)\right|,
\end{gather*}
where $\Phi$ and $\phi$ are the standard normal distribution function and
density. A \emph{two-point} (or \emph{centered Bernoulli}) summand has the
form $X_i=h_i(B_i-p_i)$, where $h_i>0$ is its support diameter and
$B_i\sim\Bern(p_i)$ ($0\le p_i\le1$); the boundary values $p_i\in\{0,1\}$
yield a deterministic zero summand. For two-point summands,
\begin{equation}\label{eq:binary-moments}
 V=\sum_{i=1}^n h_i^2p_i(1-p_i),\qquad
 R=\sum_{i=1}^n h_i^3p_i(1-p_i)\bigl\{p_i^2+(1-p_i)^2\bigr\}.
\end{equation}

By \citet[pp.~160--163]{Esseen1956}, $n$ independent copies of a two-point
summand with upper atom probability $\pe=(4-\sqrt{10})/2$ achieve
$\lim_{n\to\infty}\mathcal C=\ce:=(3+\sqrt{10})/(6\sqrt{2\pi})\approx0.409732$.
For finite $m\ge1$, \citet[Theorem~1]{Schulz2016} proved that the $m$-summand
i.i.d.\ binomial maximum $B_m:=\max_{0<p<1}\mathcal C(B_1-p,\ldots,B_m-p)$
satisfies $B_m<\ce$, and a companion paper \citep[Theorem~1.2(ii)]{BellFinite2026}
extended the upper bound $\mathcal C\le\ce$ to all finite common-diameter
two-point arrays.

Whereas \citet{BellFinite2026} bounds the ratio $\mathcal C$ from above by
$\ce$, the present paper characterizes the \emph{maximizers} of $\mathcal C$
at a fixed number of summands $n$, analyzes how $\mathcal C$ behaves under
perturbations of the probabilities and support diameters, and explains why
setting a bounded number of coordinates to zero can strictly increase the
finite-$n$ binomial maximum.

\subsection{Main results}

Let $B_m$ denote the maximal ratio for $m$ identically distributed active
two-point summands, and let $C_n^{\mathrm{com}}$ denote the maximum of
$\mathcal C(X_1,\ldots,X_n)$ over arrays of at most $n$ two-point summands
sharing a common support diameter (permitting zero summands). For $n=1$,
\begin{equation}\label{eq:c1}
 B_1=C_1=\max_{0<p<1}
 \frac{\sqrt{p(1-p)}}{p^2+(1-p)^2}
       \left\{\Phi\!\left(\sqrt{\frac p{1-p}}\right)-p\right\}
 =0.370351523474378920\ldots,
\end{equation}
which is also the unrestricted one-summand Berry--Esseen constant
\citep[Theorem~1]{BentkusKirsa1989}.

Our first theorem classifies all common-diameter maximizers at every finite
$n\ge1$.

\begin{theorem}[Classification of common-diameter maximizers]\label{thm:common}
For every integer $n\ge1$,
\[
 C_n^{\mathrm{com}}=\max_{1\le m\le n}B_m.
\]
Moreover, every $n$-summand common-diameter array attaining
$C_n^{\mathrm{com}}$ consists of $m\in\{1,\ldots,n\}$ identically distributed
nondegenerate two-point summands and $n-m$ zero summands.
\end{theorem}

When $n=2$, we remove the common-diameter hypothesis completely and determine
the exact maximum over all pairs of two-point summands with arbitrary
diameters.

\begin{theorem}[Two two-point summands with arbitrary diameters]\label{thm:two}
Over arrays of at most two independent two-point summands with arbitrary
support diameters, the maximum of $\mathcal C(X_1,X_2)$ is $C_1$, and every
attaining array has exactly one nonconstant summand. If both summands are
required to have positive variance, the supremum remains $C_1$ and is not
attained.
\end{theorem}

Our third theorem quantifies the strict improvement obtained when a
two-point array with large total count variance $A=\sum_{i=1}^np_i(1-p_i)$
and small diameter spread is replaced by the common-diameter array with the
same probabilities $p_1,\ldots,p_n$.

\begin{theorem}[Strict improvement under common-diameter replacement]\label{thm:spread}
Fix $c_0>0$. There exist constants $A_0,\sigma_0,\eta_0,\kappa>0$, depending
only on $c_0$, with the following property. Let $X_i=g_i(B_i-p_i)$
($1\le i\le n$) be independent with $g_i>0$ and $0<p_i<1$, and set
\begin{gather*}
 a_i=p_i(1-p_i),\qquad A=\sum_{i=1}^n a_i,\qquad
 h_0=\frac{\sum_{i=1}^n a_i g_i}{A},\\
 e_i=\frac{g_i}{h_0}-1,
 \qquad\sigma^2=\sum_{i=1}^n a_i e_i^2,\qquad\eta=\max_{1\le i\le n}|e_i|.
\end{gather*}
If $A\ge A_0$, $0<\sigma\le\sigma_0$, $\eta\le\eta_0$, and
$\mathcal C(X_1,\ldots,X_n)\ge c_0$, then
\[
 \mathcal C(B_1-p_1,\ldots,B_n-p_n)-\mathcal C(X_1,\ldots,X_n)
                 \ge\kappa\frac{\sigma^2}{\sigma+\eta}>0.
\]
If in addition all probabilities $p_i$ lie in a fixed compact subset
$J\subset(0,1)$, then the improvement is at least $\kappa_J\sigma$ for a
constant $\kappa_J>0$ depending only on $c_0$ and $J$. Furthermore, the
exponent $\sigma^2/(\sigma+\eta)$ cannot be improved to $\sigma$ when rare
probabilities $p_i\to0$ are permitted.
\end{theorem}

Finally, in Section~\ref{sec:iid-phase} (Theorem~\ref{thm:iid-padding}) we
determine the asymptotic behavior of the padded binomial maximum
$P_N=\max_{1\le m\le N}B_m$: for all sufficiently large $N$, $P_N$ has a
unique maximizing active count $m_N$, the number of zero summands
$j_N=N-m_N$ belongs to $\{0,1,2,3,4\}$ with explicit limiting frequencies,
and the gaps between successive record active counts belong to $\{2,3,5\}$.

\subsection{Open questions and overview of the methods}

The results of this paper motivate two open structural conjectures for fixed
$n\ge1$:
\begin{enumerate}[label=(\arabic*),leftmargin=2em]
\item \emph{Global optimality of common diameters among two-point arrays:}
Theorem~\ref{thm:two} establishes that unequal diameters cannot beat a common
diameter for $n=2$, and Theorem~\ref{thm:spread} establishes this locally for
large count variance $A$ and small diameter spread. However, as shown in the
companion numerical note, unequal diameters with rational ratios such as
$(1,1,1,2)$ at $n=4$ can create multi-summand atom collisions that are strict
local maxima of $\mathcal C$ (though with ratio $\approx0.2825<P_4\approx0.4009$).
Proving globally that $\sup\mathcal C(X_1,\ldots,X_n)=P_n$ across all
$n$-summand two-point arrays with arbitrary diameters for $n\ge3$ remains open.
\item \emph{Reduction from general summands to two-point summands:}
Combining Theorem~\ref{thm:common} with the small-Lyapunov phase-collapse and
lattice-concavity theorems of the companion structure paper proves that every
unrestricted $n$-summand global maximizer with Lyapunov ratio below an
absolute threshold $\varepsilon_0>0$ is a padded i.i.d.\ two-point array.
Extending this classification to all $n\ge1$ without a small-Lyapunov
restriction would imply $\Cind=\ce$.
\end{enumerate}

For every finite-support sum $S$, we use the one-sided branch decomposition
\begin{equation}\label{eq:branches}
 \sup_{y\in\R}\left|F(y)-\Phi\!\left(\frac y{\sqrt V}\right)\right|
 =\sup_{y\in\R}\max\left\{F(y)-\Phi\!\left(\frac y{\sqrt V}\right),\,
               \Phi\!\left(\frac y{\sqrt V}\right)-F^-(y)\right\},
\end{equation}
where $F(y)=\Pp(S\le y)$ and $F^-(y)=\Pp(S<y)$. Reflection $X_i\mapsto-X_i$
interchanges the two branches, and between consecutive atoms the closed
discrepancy $F(y)-\Phi(y/\sqrt V)$ is strictly decreasing, so its maximum is
attained at a support atom.

The paper is organized as follows:
\begin{itemize}[leftmargin=2em]
\item \textbf{Section~\ref{sec:common} (Proof of Theorem~\ref{thm:common}):}
Fixing the mean $\mu$, variance $V$, and closed integer threshold $k$ of a
common-diameter maximizer, we maximize $\sum_ip_i^3$ on the connected
component $K$ of the zero-deficit set $Z=\{p:H_C(p)=0\}$. Varying three
probabilities along their elementary-symmetric $(e_1,e_2)$ fiber and using a
uniform third-difference bound (Lemma~\ref{lem:third-difference}) reduces the
cubic maximizer on $K$ to $\Bin(M,p)+\Bern(x)$, where a certified curvature
inequality (Lemma~\ref{lem:certified-curvature}) rules out $p\ne x$ and proves
that the i.i.d.\ point is isolated in $Z$, forcing $K$ to be a singleton.
\item \textbf{Section~\ref{sec:two} (Proof of Theorem~\ref{thm:two}):}
For two two-point summands with diameters $1$ and $t\in(0,1)$, variational
contact bounds confine any stationary point with $\mathcal C\ge37/100$ to a
compact variance domain $V>1369/3319684$, where a 37,287-leaf certified
interval partition rules out all three interior threshold branches.
\item \textbf{Section~\ref{sec:iid-phase} (Binomial phase and zero padding):}
We derive a uniform seventh-order Edgeworth--Euler--Maclaurin expansion for
$D_{m,k}(v)=F_{m,v}(k)-\Phi(-cd)$, compute the $O(m^{-2})$ and $O(m^{-3})$
phase expansions of $B_m$, determine the five-parabola lower envelope on
$[-1/2,1/2]$, and use a field-norm lower bound in $\mathbb Q(\sqrt{10})$ to
control the $O(N^{-1})$ boundary drift.
\item \textbf{Section~\ref{sec:spread} (Proof of Theorem~\ref{thm:spread}):}
Using uniform Poisson-binomial local and jittered Edgeworth expansions with
no individual probability floor, we bound the unconditioned weighted
perturbation moments $N_j(k)=\E[T^j\ind_{\{K=k\}}]$, localize any threshold
with $\mathcal C\ge c_0$ to the central bulk, establish the conditional $L^1$
lower bound $\E(|T|\mid K=k)\ge c\ell/\sqrt A$ by a perfect-matching coupling
of two rejective-sampling copies, and construct the rare-coordinate example
proving optimality of the rate $\sigma^2/(\sigma+\eta)$.
\end{itemize}

\section{All finite common-diameter maximizers}\label{sec:common}
Scale the common diameter to one and write
\[
 \mu=\sum_i p_i,\quad V=\sum_i p_i(1-p_i),\quad
 R=\sum_i\{p_i-3p_i^2+4p_i^3-2p_i^4\},\quad
 F_k=\Pp\{\textstyle\sum_i B_i\le k\}.
\]
Since $R\ge V/2$, the ratio is at most $2\sqrt V$.
Its extension by zero at $V=0$ is continuous on the compact probability
cube. Thus $C=C_n^{\mathrm{com}}$ is attained.
At any selected closed integer threshold, define
\begin{equation}\label{eq:deficit}
 H_C(p)=\frac{CR}{V^{3/2}}+
             \Phi((k-\mu)/\sqrt V)-F_k.
\end{equation}
The global definition of $C$ gives $H_C\ge0$ everywhere with $V>0$,
and equality at every maximizing array at its selected threshold.

The probability replacements below have classical precedents.
Hoeffding's pair identity preserves the sum of two probabilities
\cite[equation~(10), pp.~714--716]{Hoeffding1956}.
Mattner's few-moment extremum theorem and next-power-sum argument
\cite[Theorems~1.1--1.2 and Lemma~2.1]{Mattner2022} use related fibers
for separately affine symmetric functions. The present deficit contains
a cubic-moment denominator. We establish its required affinity directly.

\subsection{A uniform third-difference estimate}
For the mass function $c_j$ of a Poisson-binomial sum, extended by zero,
put $\Delta^3c_j=c_j-3c_{j-1}+3c_{j-2}-c_{j-3}$.

\begin{lemma}\label{lem:third-difference}
If that sum has variance $W$, then
\[
                \sup_j|\Delta^3c_j|<\frac{16/5}{(1+W)^{3/2}}.
\]
\end{lemma}
\begin{proof}
For $W\le1/4$, choose the more likely outcome of each bit as a base
configuration. Let $J$ count departures, $\epsilon_i=\min(p_i,1-p_i)$,
$E=\sum\epsilon_i$, and $Q=\sum\epsilon_i/(1-\epsilon_i)$.
Then $E\le2W\le1/2$, $W\le Q\le1$,
\[
 \Pp(J=0)\ge e^{-Q},\qquad
 \Pp(J=1)=Q\Pp(J=0)\ge We^{-W}.
\]
At the two indices where the base state contributes magnitude three,
the one-departure terms have the opposite sign and magnitude at least
one. The base probability is at least $1/2$, so the sign cannot reverse.
Thus the absolute third difference is at most $3-4We^{-W}$ there.
At all other indices it is at most $1+4W$.
For the central bound use $e^{-W}\ge1-W$ and
\[
 (1+W)^{-3/2}\ge h(W):=
       1-3W/2+15W^2/8-35W^3/16.
\]
The polynomial $(16/5)h(W)-(3-4W+4W^2)$ is decreasing on $[0,1/4]$:
its derivative is bounded above by
$-1/2+13W/4-315W^2/16<0$.
At $W=1/4$ it equals $1/64$.
The outer estimate is at most two, below $(16/5)(5/4)^{-3/2}$.

For $W\ge1/4$, Fourier inversion gives
\[
 |\Delta^3c_j|\le I(W):=
 \frac{16}{\pi}\int_0^{\pi/2}\sin^3u\,e^{-2W\sin^2u}\,du.
\]
Write $J_r=\int_0^{\pi/2}\sin^ru\,e^{-2W\sin^2u}\,du$.
Integration of $(\sin^4u\cos u\,e^{-2W\sin^2u})'$ gives
$(5+4W)J_5=4J_3+4WJ_7$. Consequently
\[
 \frac{d}{dW}\log\{I(W)(1+W)^{3/2}\}
 \le-\frac8{5+4W}+\frac3{2(1+W)}<0.
\]
At $W=1/4$, the exponential upper Taylor polynomial through degree six
and the exact odd sine integrals give
\[
 I(1/4)(5/4)^{3/2}
 \le\frac{10\sqrt5}{\pi}\,
       \frac{911264}{2027025}<\frac{16}{5}.
\]
The last strict inequality follows by squaring and using
$\pi>6283/2000$; it is one of the rational scalar controls in
the computational supplement.
\end{proof}

\subsection{Reducing the number of distinct probabilities}
Fix a maximizing array and its mean, variance and closed threshold.
Consider the compact zero set
\[
 Z=\{p\in[0,1]^n:\ \sum p_i=\mu,\
                  \sum p_i(1-p_i)=V,\ H_C(p)=0\}.
\]
Let $K$ be its connected component containing the array, and maximize
$\sum p_i^3$ on $K$. For three selected probabilities, let $e_1,e_2,e_3$
be their elementary symmetric functions. Their generating polynomial is
\[
 1+e_1(w-1)+e_2(w-1)^2+e_3(w-1)^3.
\]
Holding $e_1,e_2$ fixed preserves their mean and variance. Their
contribution to $R$ is
\begin{align*}
 R_3={}&e_1-3e_1^2+4e_1^3-2e_1^4
 +(6-12e_1+8e_1^2)e_2-4e_2^2+4(3-2e_1)e_3.
\end{align*}
Both $F_k$ and $R$ are affine in $e_3$ on this fiber. If its three
roots are distinct and interior, $e_3$ can move in both directions.
A nonzero slope contradicts $H_C\ge0$. A zero slope gives a continuous
path in $Z$, hence in $K$, on which $\sum p_i^3=e_1^3-3e_1e_2+3e_3$
increases. The chosen point therefore has at most two distinct active
probabilities.

Suppose they are $a<b$, each occurring at least twice. The triples
$(a,a,b)$ and $(a,b,b)$ are opposite endpoints of their three-root
fibers. If $d_j$ is the mass function of the other coordinates, the
fiber slope is
\[
 L=\frac{4C(3-2e_1)}{V^{3/2}}+\Delta^2d_k.
\]
The two endpoint conditions give $L_{aab}\le0\le L_{abb}$.
Delete two bits of each type, let $c_j$ be the remaining mass function,
and let $W$ be its variance. Direct subtraction gives
\[
 L_{abb}-L_{aab}
       =(b-a)\left\{\Delta^3c_k-\frac{8C}{V^{3/2}}\right\}.
\]
Since $V\le W+1$, Lemma~\ref{lem:third-difference} contradicts these
signs whenever $C\ge2/5$. Thus one of the two multiplicities is one.
At $n=3$ this is automatic. At $n=4$, the only remaining repeated
pattern is $(2,2)$; its deleted law is deterministic, so
$|\Delta^3c_k|\le3$, whereas $8C/V^{3/2}\ge8C>3$ when
$C\ge199/500$.

\subsection{Strict curvature of the remaining fiber}
The remaining unequal pattern is $\Bin(M,p)+\Bern(x)$, $p\ne x$.
We record its exact curvature and the finite inequality used to exclude it.
Let $q=1-p$, $m=M-1$, $d_j=\Pp(\Bin(m,p)=j)$, $W=mpq$, and
\[
 V=x(1-x)+Mpq,\quad z=\frac{k-x-Mp}{\sqrt V},\quad
 \xi=d_{k-1}-d_k,\quad
 S_2=(1-2x)^2+(1-2x)(1-2p)+(1-2p)^2.
\]
Preserve the sum of the exceptional mass and one mass $p$, write them
as $\bar p+d,\bar p-d$, and set $u=d^2$. Along this path,
\[
 F_k'=\xi,\qquad V'=-2,\qquad R'=-2S_2,\qquad R''=-8.
\]
At stationarity of $H_C$, differentiating its smooth terms and
eliminating $R$ gives
\begin{equation}\label{eq:fiber-curvature}
 V^2H_C''=5V\xi-C\left(8\sqrt V+\frac{2S_2}{\sqrt V}\right)
                       -z\phi(z)(2+z^2).
\end{equation}
A minimum requires this expression to be nonnegative.

\begin{lemma}[The certified curvature inequality]\label{lem:certified-curvature}
Put $c_*=199/500$. For
\[
 0<p,x<1,\quad pq\le W\le3,\quad0\le f\le1,\quad
 A_0=W-(1-f)q^2\ge0,
\]
define
\[
 B_0=p^2+(1-f)(W+2pq),\quad V=x(1-x)+W+pq,
 \quad z=\frac{f(W+pq)/(1-qf)-x}{\sqrt V}.
\]
Then
\begin{equation}\label{eq:compact-curvature}
 \frac{2c_*}{\sqrt V}\{4W+2+(1-2x)(1-2p)\}
 +z\phi(z)(2+z^2)
 -\frac{5VfB_0}{(2-f)B_0+A_0}>0.
\end{equation}
Also every two-bit unit-diameter array has ratio less than $c_*$.
\end{lemma}

The lemma is a finite arithmetic premise. Its compact interval cover
at 192 Arb bits \cite{Johansson2017} processed 550563 boxes for
\eqref{eq:compact-curvature} and 13641 boxes for the two-bit assertion,
with no unresolved box. The archived records bind the evaluated
formulas and the declared domains. Below we explain the covering
charts and the analytic complements; rebuilding those interval covers
is distinct from checking the saved hashes and coverage metadata.

For $V\le1/5$, the first term in \eqref{eq:compact-curvature} is at
least $8c_*/\sqrt3$, the normal term is at least $-4/(e\sqrt\pi)$,
and the final term is at most one. Rational controls give
\begin{equation}\label{eq:scalar-margins}
 \frac{8c_*}{\sqrt3}-\frac4{e\sqrt\pi}-1>\frac1{247},\quad
 8c_*\sqrt3-\frac{10}{\pi}-\frac4{e\sqrt\pi}>1,\quad
 8c_*\sqrt{7/2}-\frac{35}{3\pi}-\frac4{e\sqrt\pi}>\frac{44}{45}.
\end{equation}
For $V\ge1/5$ use
$z\phi(z)(2+z^2)\ge\sqrt{2/\pi}\min(z,0)$.
The following maps from $[0,1]^4$, with fourth variable $x$, cover
the entire remaining domain:
\begin{align*}
 &p=(1+s_p)/2,\quad W=pq+s_W(3-pq),\quad f=s_f;\\
 &p=s_p/2,\quad W=q^2+s_W(3-q^2),\quad f=s_f;\\
 &p=s_p/2,\quad W=pq+s_W(q^2-pq),\quad
                    f=1-(1-s_f)W/q^2.
\end{align*}
The first covers $p\ge1/2$. The last two split $p\le1/2$ at $W=q^2$;
below that interface the last map imposes exactly $f\ge1-W/q^2$.
Interval division is used only with a strictly positive denominator.
Otherwise the universal neighbor-fraction bound one, or $y\ge0$ at a
singular endpoint, supplies a valid lower bound. A box is accepted only
when its outward lower endpoint is positive. The two-bit cover includes
the entire probability square and all three closed integer thresholds.
The zero-variance corners use $\mathcal C\le2\sqrt V$ before division.

To apply the lemma to \eqref{eq:fiber-curvature}, first suppose $M\ge2$.
If $\xi\le0$, then $\max(V,S_2)\ge1/3$ and
\[
 C(8\sqrt V+2S_2/\sqrt V)+z\phi(z)(2+z^2)
 \ge8c_*/\sqrt3-4/(e\sqrt\pi)>0.
\]
Thus $H_C''<0$. If $\xi>0$, binomial recurrence gives
\[
 y=k-(m+1)p>0,\quad f=y/(kq)\in(0,1],\quad
 d_k=(1-f)d_{k-1},\quad d_{k-2}=r_0d_{k-1},
\]
where
\[
 r_0=A_0/B_0,\qquad y=f(W+pq)/(1-qf).
\]
The sum of these three masses is at most one, so
$\xi\le f/(2-f+r_0)$. For $W\le3$, Lemma~\ref{lem:certified-curvature}
therefore makes \eqref{eq:fiber-curvature} negative. The coefficient
of $C$ is positive, so the same conclusion holds for every actual extremal ratio
$C\ge c_*$.
For $W\ge3$, Fourier inversion gives
\[
 |\xi|\le\frac4\pi\int_0^{\pi/2}\sin u\,e^{-2W\sin^2u}\,du
                        <\frac2{\pi W}.
\]
The negative of the right side of \eqref{eq:fiber-curvature} is at least
$8c_*\sqrt V-4/(e\sqrt\pi)-10V/(\pi W)$.
It is concave in $V\in[W,W+1/2]$. Both endpoint values increase with
$W$, and \eqref{eq:scalar-margins} gives positive bounds at $W=3$.
For $M=1$ the two-bit assertion of the lemma excludes the pattern at
price $C\ge c_*$.

\subsection{From one equal-mass point to every maximizer}
The cubic maximizer on $K$ is now an IID group plus deterministic bits.
The number $m$ of active bits is at least three, by
Lemma~\ref{lem:certified-curvature}. We prove this point is isolated in $Z$.
For a pair in the IID group, use $p+d,p-d$ and $u=d^2$.
Nonnegativity gives $\Lambda:=\partial_uH_C(0)\ge0$.
If $\Lambda=0$, \eqref{eq:fiber-curvature} applies at the diagonal
with deleted law $\Bin(m-2,p)$. The compact lemma or its analytic
complement makes $\partial_u^2H_C(0)<0$, contradicting $H_C\ge0$
for small positive $u$. Hence $\Lambda>0$.
Permutation symmetry makes the exchange quadratic form
$(\Lambda/2)\|y_{\mathrm{cen}}\|^2$ on the mean-zero directions
of the active probabilities.

For a nearby point in $Z$, let $a_i\ge0$ activate zero probabilities,
$b_j\ge0$ deactivate probability-one coordinates, and $y_l$ change
the $m$ active probabilities. Put $a=\sum a_i$, $b=\sum b_j$, $t=a+b$.
The two constraints give exactly
\[
 \sum_l y_l=b-a,\qquad
 \sum_l y_l^2=2pa+2(1-p)b-\sum_i a_i^2-\sum_j b_j^2.
\]
Thus $\|y_{\mathrm{cen}}\|^2\ge c_pt-O(t^2)$, with $c_p>0$.
The endpoint linear coefficients of $H_C$ are nonnegative: activate
one endpoint and compensate uniformly in the active group to preserve
the mean, then apply $H_C\ge0$ on that feasible path.
The Taylor expansion on the fixed mean/variance fiber consequently has
nonnegative endpoint linear terms, the positive exchange term, and
remaining terms $O(t^{3/2})$. The uniform active direction is $O(t)$,
the exchange direction is $O(\sqrt t)$, so these remainder orders also
include their mixed terms. For small $t>0$ this gives $H_C>0$.
If $t=0$, the two constraints force $y=0$. The IID point is isolated.
Its connected component $K$ is therefore a singleton, and the original
maximizer was the same IID group plus zeros.

The required lower prices follow from two exact admissible examples:
\[
 \begin{array}{c|c|c|c}
 \text{active count}&p&\text{closed raw threshold}&\text{ratio lower bound}\\\hline
 3&43/100&1&0.3981444741>199/500\\
 5&9/20&2&0.4009890598>2/5.
 \end{array}
\]
Their binomial probabilities are rational. Rational square-root bounds,
Machin bounds for $\pi$, and the integrated Gaussian Taylor polynomial
$x-x^3/6+x^5/40-x^7/336$ give the displayed strict lower bounds.
Padding preserves them. The three-bit example handles $n=3,4$;
the five-bit example supplies the $2/5$ hypothesis for the repeated-type
argument at every $n\ge5$.

\subsection{The remaining two-coordinate case}
For two probabilities $p,x$, set $s=p+x$, $w=px$ and
\[
 V=s-s^2+2w,\quad
 R=s-3s^2+4s^3-2s^4+(6-12s+8s^2)w-4w^2.
\]
At rank $k=0,1,2$, the closed CDFs are $1-s+w$, $1-w$, $1$.
Let $z=(k-s)/\sqrt V$, $D=F_k-\Phi(z)$ and $C=DV^{3/2}/R$.
Off the diagonal the map $(p,x)\mapsto(s,w)$ is invertible, so an
interior local maximum requires both derivatives to vanish.
With $F_w=(1,-1,0)_k$, $F_s=(-1,0,0)_k$,
$V_s=1-2s$, $R_w=6-12s+8s^2-8w$ and
$R_s=(1-2s)^3+(-12+16s)w$, their positive-denominator numerators are
\begin{align*}
 G_w={}&R(VF_w+z\phi(z)+3D)-DVR_w,\\
 G_s={}&2VRF_s+2\sqrt V R\phi(z)+Rz\phi(z)V_s
                           +3DRV_s-2DVR_s.
\end{align*}
At $G_w=0$,
\[
 \frac{C_{ww}}C=\frac{N}{DV^2R},\quad
 N=R\{6F_wV+z\phi(z)(3+z^2)+3D\}+8DV^2.
\]
The finite critical certificate covers each owner square by
$C<37/100$, $0\notin G_w$, $0\notin G_s$, or $N>0$.
A separate diagonal cover uses only the ratio bound, since the
$(s,w)$ chart is singular there. These covers processed respectively
4683 and 195 boxes, with no unresolved box. They exclude every
two-active-bit common-diameter maximizer at ratio at least $37/100$.
Boundary derivatives are not assumed to vanish; the cover bounds their
limits only to enclose the open stationary domain.
The exact one-bit example $p=2/5$ gives a ratio greater than $37/100$.
Compact attainment therefore leaves exactly one active bit at $n=2$.
The $n=1$ case is immediate. This completes
Theorem~\ref{thm:common}.

\section{Two arbitrary binary diameters}\label{sec:two}
We first justify attainment and separate the boundary cases from the
open stationary domain. Normalize $V=1$ and take a maximizing sequence
with ratio at least $c>0$. For two coordinates, $R\ge1/\sqrt2$,
whereas the discrepancy is at most one, so $R\le1/c$.
The squares are uniformly integrable, since
\[
 \sum_i\E[X_i^2;|X_i|>L]\le\frac1{cL}.
\]
After reflection, selected positive closed thresholds stay bounded:
Chebyshev controls thresholds tending to $-\infty$, and $1-\Phi(y)$
controls those tending to $+\infty$. Coordinatewise weak limits preserve
centering and total variance. A weak limit of two-point laws has at most
two points. Closed-set Portmanteau at the moving thresholds and lower
semicontinuity of $R$ give a limiting ratio at least the supremum;
equality follows from its definition. Thus a binary-class maximizer exists.

\subsection{The top threshold and equal-diameter boundary}
For a centered binary variable on $\{-b,a\}$, $a,b>0$,
\[
 v=ab,\qquad \rho=\frac{ab(a^2+b^2)}{a+b}
                 \ge2(\sqrt2-1)\frac{v^2}{a}.
\]
The inequality is the minimum of $(1+x^2)/(1+x)$ for $x=a/b>0$.
For two such coordinates with total upper endpoint $L$, Cauchy's
inequality gives $R\ge2(\sqrt2-1)V^2/L$.
At the top closed threshold the CDF is one, so its ratio is at most
\[
 \frac{z\{1-\Phi(z)\}}{2(\sqrt2-1)}
 \le\frac{e^{-1/2}}{4(\sqrt2-1)}<\frac{37}{100},
 \qquad z=L/\sqrt V.
\]
Here $1-\Phi(z)\le e^{-z^2/2}/2$ follows by translating its Gaussian
integral. For the final rational comparison,
$\sqrt e>41/25$, $\sqrt2-1>207/500$ give
$12500/33948<37/100$.
Reflection treats the negative strict bottom threshold.
The common-diameter two-active-bit case was excluded in
Section~\ref{sec:common}. A zero coordinate leaves exactly the
one-coordinate problem, whose ratio exceeds $37/100$.

\subsection{Open unequal-diameter stationary points}
Scale the larger diameter to one and write the other as $t\in(0,1)$.
Let their upper probabilities be $p,q\in(0,1)$, and put
\[
 a=p(1-p),\quad b=q(1-q),\quad
 r_p=p^2+(1-p)^2,\quad r_q=q^2+(1-q)^2,
 \quad V=a+t^2b,\quad R=ar_p+t^3br_q.
\]
The raw support is $0,t,1,1+t$. The three remaining positive closed
thresholds have the data
\[
 \begin{array}{c|c|c}
 \text{raw threshold}&F&y\ \text{(centered threshold)}\\\hline
 0&(1-p)(1-q)&-p-tq\\
 t&1-p&-p+t(1-q)\\
 1&1-pq&1-p-tq.
 \end{array}
\]
Set $z=y/\sqrt V$, $D=F-\Phi(z)$ and $C=DV^{3/2}/R$.
At a global maximum, a selected active branch has a local maximum
even if other thresholds tie. Its three derivatives therefore vanish.
Define
\[
 \alpha=\phi(z)/\sqrt V,\quad
 A=\frac{3D+z\phi(z)}{2V},\quad\lambda=D/R.
\]
After removing the common positive factor $V^{3/2}/R$, the two
probability derivatives are
\begin{align*}
 g_p&=F_p+\alpha+A(1-2p)-\lambda(1-2p)^3,\\
 g_q&=F_q+\alpha t+At^2(1-2q)-\lambda t^3(1-2q)^3.
\end{align*}
Divide the diameter derivative by $q$ for thresholds $0,1$ and by
$1-q$ for threshold $t$, then multiply by $VR\sqrt V$. Its numerator is
\begin{equation}\label{eq:gap-gradient}
 G=\phi(z)RB+3Dwt a\sqrt V\,(r_p-tr_q),
\end{equation}
where the pairs $(w,B)$ for these three rows are, respectively,
\[
 \bigl(1-q,\ p[(1-p)-t(1-q)]\bigr),\quad
 \bigl(q,-p[(1-p)+qt]\bigr),\quad
 \bigl(1-q,(1-p)[p+t(1-q)]\bigr).
\]
The extra factors $p$ and $1-p$ in the first and last rows are
strictly positive in the open chart and can also be removed.
These factorizations retain cancellation near vanishing probabilities.

\subsection{A compact domain for the necessary equations}
The probability equations also bound the normalized support points.
For one centered binary coordinate with law $P_i$, put
$s_i=\E X_i^2$, $b_i=\E|X_i|^3$, $M_i=\E X_i|X_i|$ and let
$G_i$ be the deleted CDF. In total-variance-one units define
\[
 Q_i(u)=G_i(y-u)-F(y)+\phi(y)u+A(u^2-s_i)
                 -C\{|u|^3-b_i-3M_i u\}.
\]
Its mean under $P_i$ is zero. Probability stationarity says its two
endpoint values agree, so both equal zero. This uses probability
stationarity alone, without a full-law majorant.
Since $Cb_i\le CR=D\le1$, $|M_i|\le s_i\le1$,
$\phi<2/5$ and $|A|\le13/8$, every active endpoint satisfies
\[
 C|u|^3\le4+(2/5+3C)|u|+(13/8)u^2.
\]
For $|u|\ge3$, dividing the right side by $u^2$ bounds it by
$C+793/360$. In particular $|u|<3+8/C$.
At $C\ge37/100$, the larger normalized diameter $1/\sqrt V$
is less than twice $911/37$. Hence
\begin{equation}\label{eq:two-compact-domain}
 V>\frac{1369}{3319684}>\frac1{2500},\qquad
 R\ge V^{3/2}/\sqrt2>V/75.
\end{equation}

The arithmetic certificate covers the three closed cubes
$(p,q,t)\in[0,1]^3$, conditional on \eqref{eq:two-compact-domain}
and $C\ge37/100$. Each box is excluded by an empty conditional domain,
a strict ratio upper bound, or a nonzero enclosure of a required
probability or factored diameter derivative. One useful ratio bound is
\[
 C\le D\sqrt{a/r_p^2+b/r_q^2},
\]
which follows from H\"older's inequality applied to $V,R$.
There are 37287 accepted source leaves in the complete partition.
An independently written, cancellation-preserving replay checks every
source condition, using 37421 descendant leaves and no unresolved
or unstarted region. The input bindings and exact binary forests are
part of the public certificate. Boundary boxes enclose limits of open
stationary points; they do not impose stationarity on actual one-bit
boundary maximizers.

Consequently no competitive unequal-diameter two-active-bit maximum
can satisfy its necessary equations. The preceding endpoint and
common-diameter arguments cover all remaining two-active-bit cases.
Attainment leaves exactly one active coordinate, proving
Theorem~\ref{thm:two}.
To require both variances positive, add a centered symmetric binary
coordinate of diameter tending to zero. Choose the closed threshold
just above the winning one-coordinate atom. Its CDF, variance and
third-moment denominator converge to the one-coordinate values.
Thus the supremum remains $C_1$, and the classification prevents its
attainment with both coordinates active.

\section{Binomial phase and the optimal number of zero coordinates}
\label{sec:iid-phase}

This section concerns the binomial optima $B_m$ and their padded envelope
$P_N:=\max_{1\le m\le N}B_m$. It determines the eventual record pattern
within this family. Theorem~\ref{thm:common} transfers these conclusions
to common-diameter maximizers; no reduction from arbitrary independent
laws is assumed here.

Put $s=\sqrt{10}$ and write
\begin{equation}\label{eq:phase-constants}
\begin{gathered}
 p=(4-s)/2,\quad q=1-p,\quad t=s-3,\quad u=p-t,
 \quad c_*=(20+s)/72,\\
 Q_*=\frac{\phz(867+287s)}{25920},\qquad
 \alpha=\frac{\phz(60+19s)}{30},\qquad
 b=\frac{1010-397s}{12960},\\
 \delta=c_*-b=\frac{2590+577s}{12960},\qquad
 \eta=Q_*/\alpha=\frac{2510-747s}{8640},\qquad
 \omega=\frac{7773-2447s}{8640}.
\end{gathered}
\end{equation}
The letters $p,q,t$ in this section without an argument denote these
fixed limiting values. For a variable Bernoulli parameter $v$, use
$a(v)=v(1-v)$, $r(v)=v^2+(1-v)^2$ and $\tau(v)=1-2v$.
For a real number $x$, write $\langle x\rangle=x-\operatorname{round}(x)$,
choosing either convention at half-integers.

\begin{theorem}[Eventual binomial padding]\label{thm:iid-padding}
At every sufficiently large horizon $N$, $P_N$ has a unique maximizing
core size $m_N$, and $j_N=N-m_N$ belongs to $\{0,1,2,3,4\}$.
The limiting frequencies of these five zero counts are respectively
\[
 p-\omega,\quad p-\omega,\quad t,\quad\omega,\quad\omega.
\]
The gaps between successive record core sizes eventually belong exactly
to $\{2,3,5\}$: no other gap occurs sufficiently late, and each of
these three occurs infinitely often. No effective first horizon is asserted.
\end{theorem}

\subsection{A uniform expansion at a fixed integer threshold}

Let $F_{m,v}$ be the distribution function of $\Bin(m,v)$. Fix an integer
$k$ and set $c=mv-k$, $d=(ma(v))^{-1/2}$. The positive closed-threshold
discrepancy is $D_{m,k}(v)=F_{m,v}(k)-\Phi(-cd)$.

\begin{lemma}\label{lem:phase-uniform}
On every compact $J\subset(0,1)$ and bounded range $|c|\le M$, uniformly
for $v=(k+c)/m\in J$,
\begin{equation}\label{eq:phase-cdf}
 D_{m,k}(v)=\phz\bigl[d(1/2+\tau/6)+d^3C(\tau,-c)
                         +d^5T_5(\tau,-c)\bigr]+O(d^7),
\end{equation}
with the same order for the derivative of the remainder in $c$, fixing
$m,k$. The expansion extends to any fixed order and any fixed finite
number of these branch derivatives. Here
\begin{align*}
C(\tau,\beta)={}&\frac{-1080(1+\tau)\beta^2
 -(90\tau^2+1080\tau+630)\beta
 +\tau^3-45\tau^2-189\tau-135}{4320},\\
T_5(\tau,\beta)={}&\frac{3+5\tau}{48}\beta^4
 +\frac{27+60\tau+25\tau^2}{288}\beta^3
 +\frac{15+31\tau+25\tau^2+\tau^3}{192}\beta^2\\
&+\frac{549+1320\tau+1030\tau^2+120\tau^3+5\tau^4}{23040}\beta\\
 &+\frac{189+927\tau+126\tau^2+234\tau^3+21\tau^4-25\tau^5}{193536}.
\end{align*}
\end{lemma}

\begin{proof}
The Fourier integral for a binomial mass is on $[-\pi,\pi]$. Its modulus
satisfies
\[
 |1-v+ve^{i\xi}|^m
 =(1-4a(v)\sin^2(\xi/2))^{m/2}.
\]
After $z=\xi\sqrt{ma(v)}$, the inequality
$\sin(|\xi|/2)\ge|\xi|/\pi$ bounds this by $e^{-2z^2/\pi^2}$
throughout the integration interval. On $J$, the centered logarithm has
a Taylor expansion of any prescribed order with uniformly bounded
coefficients and parameter derivatives. Split the integral at
$|z|=m^\varepsilon$, with $\varepsilon$ sufficiently small for that
order. Polynomial multiples of the Gaussian bound control the inner
Taylor remainder; the outer integral is smaller than any prescribed
power of $m^{-1}$. Taking sufficiently many terms gives a Hermite
mass expansion with uniform absolute error $O(d^{12})$, or any higher
fixed power if needed, at every integer mass index.

This statement also holds for branch derivatives. Before centering,
\[
 \frac{\partial}{\partial c}(1-v+ve^{i\xi})^m
 =(e^{i\xi}-1)(1-v+ve^{i\xi})^{m-1},
 \qquad v=(k+c)/m.
\]
Further fixed derivatives have the same Gaussian control after increasing
the retained Taylor order. The standardized mass index and spacing have
bounded branch derivatives of the needed orders on the stated ranges.

Summing at most $m+1=O(d^{-2})$ masses turns the displayed mass error
into $O(d^{10})$, including its branch derivatives. The tails beyond
the true binomial support are exponentially small. Euler--Maclaurin
summation may therefore be applied to the infinite lower lattice of the
Gaussian-polynomial terms. Its order-eight remainder is bounded by a
constant times $d^8\|f^{(8)}\|_1$ for each such term. In the integral
periodic-Bernoulli remainder, differentiating the phase gives bounded
factors of the form $1+d|z|$; Gaussian derivative moments remain
integrable. Increasing the summation order gives the corresponding
assertion for further branch derivatives.

The coefficient of $d^j$ in the density has parity $j$. Integration,
the half-cell term, and the Euler--Maclaurin derivative terms then leave
only odd powers $d,d^3,d^5,\ldots$ when evaluated at $-cd$.
The terms through $d^5$ give \eqref{eq:phase-cdf}. Their coefficients
follow by using the Bernoulli cumulant ratios
\[
 \tau,\quad1-6a,\quad\tau(1-12a),\quad
 1-30a+120a^2,\quad\tau(1-60a+360a^2)
\]
and the lattice corrections $df/2+d^2f'/12-d^4f'''/720$.
The cubic phase coefficient is also the polynomial in Schulz's
Remark~4.3, derived on p.~25 \cite{Schulz2016}.
\end{proof}

Dividing \eqref{eq:phase-cdf} by $r(v)d$ gives the threshold efficiency
\begin{equation}\label{eq:phase-efficiency}
 R_{m,+,k}(v)=E(v)-\frac{Q(v,c)}m-\frac{H(v,c)}{m^2}+O(m^{-3}),
\end{equation}
where
\[
 E(v)=\phz\frac{3+\tau(v)}{3(1+\tau(v)^2)},\quad
 Q(v,c)=-\frac{\phz C(\tau(v),-c)}{r(v)a(v)},\quad
 H(v,c)=-\frac{\phz T_5(\tau(v),-c)}{r(v)a(v)^2}.
\]
The phase parabola $Q(v,c)$ has vertex
\[
 c_0(v)=\frac{\tau(v)^2+12\tau(v)+7}{24(1+\tau(v))}
\]
and minimum
\begin{equation}\label{eq:phase-qmin}
 Q_{\min}(v)=\frac{\phz\{345+72\tau-498\tau^2-8\tau^3-23\tau^4\}}
 {4320(1+\tau)(1-\tau^4)}.
\end{equation}

\subsection{Global binomial localization and the padding envelope}

Schulz's Theorem~1 gives $B_m<\ce$ for every finite $m$.
His Proposition~6.1 bounds the full Kolmogorov distance on
$v\in[1/6,5/6]$ by the four absolute closed/left central-threshold errors
at $\lfloor mv\rfloor,\lceil mv\rceil$ and
\[
 \gamma\frac{\phz(3+|1-2v|)}{6\sqrt{mv(1-v)}},
 \qquad \gamma=\frac{2}{\sqrt{2\pi}\ce}-1=12t-1<1.
\]
His estimate on the two complementary parameter ranges is
$R_m(v)\le .3031(1+.646/\beta(v))<.4045$,
where $\beta(v)=r(v)/\sqrt{a(v)}$
\cite[pp.~29--31]{Schulz2016}.
Each $B_m$ is attained: the ratio is the maximum of finitely many
continuous threshold expressions and tends to zero as $v\to0,1$.

By reflection restrict $v\le1/2$. Central expansions and the source
bound show uniform convergence of the optimized threshold value to $E(v)$
on the retained compact parameter interval. Its unique maximum is
$\ce$ at $v=p$. Thus every maximizing parameter tends to $p$.
Near $p$, the extra $\gamma$ term and left-limit thresholds are uniformly
smaller than the positive closed central threshold: the latter has leading
coefficient $\phz(1/2+\tau/6)/r$, while the left-limit coefficient is
$\phz(1/2-\tau/6)/r$. Consequently the full winner is eventually
a positive closed floor/ceiling threshold, with $|c|\le1$.

Taking $v=(k+c_*)/m$ with $k$ nearest $mp-c_*$ gives
$\ce-Q_*/m+O(m^{-2})$. Since
$E(v)=\ce-\alpha(v-p)^2+O((v-p)^3)$, comparison with
\eqref{eq:phase-efficiency} and $Q\ge Q_{\min}$ yields
\[
 (v-p)^2\le C|v-p|/m+C/m^2,
 \qquad v-p=O(m^{-1}).
\]
The positive quadratic phase penalty then gives
$c-c_0(v)=O(m^{-1/2})$. Since $0<c_*<1$, the branch is eventually
interior. Differentiating \eqref{eq:phase-efficiency} along
$v=(k+c)/m$ and using $E'(v)=O(m^{-1})$ gives $Q_c=O(m^{-1})$,
so $c-c_0(v)=O(m^{-1})$ and
$v=(k+c_*)/m+O(m^{-2})$.

Set $e=k+c_*-mp$. Taylor expansion and completion of the square give
\begin{equation}\label{eq:phase-bm}
 B_m=\ce-\frac{Q_*}m-\frac{H_*+
 \alpha\operatorname{dist}(mp-\delta,\Z)^2}{m^2}+O(m^{-3}).
\end{equation}
Indeed the second-order branch correction is
$\alpha e^2+Q_p e+H_0$, where all derivatives are evaluated at
$(p,c_*)$, $H_0=H(p,c_*)$, $b=-Q_p/(2\alpha)$ and
$H_*=H_0-\alpha b^2$. Direct substitution in
\eqref{eq:phase-qmin} gives \eqref{eq:phase-constants}.
The formula remains uniform at the nearest-integer crossing, by taking
the better of the two branch values; the nearest branch itself need
only be selected outside an $O(m^{-1})$ half-integer neighborhood.

Choose the most recent maximizing core $m_N$ in a tie. Strict finite
$B_m<\ce$ and $B_m\to\ce$ imply $m_N\to\infty$.
The leading correction in \eqref{eq:phase-bm}, compared with $B_N$,
then gives $m_N/N\to1$, and its $O(m^{-2})$ bounds give
$j_N=N-m_N=O(1)$. With
$\theta_N=\operatorname{round}(Np-\delta)-(Np-\delta)$, for bounded $j$,
\[
 N^2(B_{N-j}-B_N)=\alpha\{\theta_N^2-
 \operatorname{dist}(\theta_N+jp,\Z)^2-\eta j\}+O(N^{-1}).
\]
The limiting policy minimizes
$f_j(\theta)=\eta j+\operatorname{dist}(\theta+jp,\Z)^2$.
Define
\begin{equation}\label{eq:phase-boundaries}
\begin{aligned}
 T_{13}&=-(p+u)/2+\eta/t,& T_{30}&=-u/2-3\eta/(2u),\\
 T_{02}&=t/2+\eta/t,&T_{24}&=3t/2+\eta/t,\\
 T_{41}&=(q+2t)/2-3\eta/(2u).
\end{aligned}
\end{equation}
In increasing order on $[-1/2,1/2]$ the winning zero counts are
\begin{center}
\begin{tabular}{cc}
\toprule
Phase interval&Zero count\\
\midrule
$[-1/2,T_{13}]$&1\\
$[T_{13},T_{30}]$&3\\
$[T_{30},T_{02}]$&0\\
$[T_{02},T_{24}]$&2\\
$[T_{24},T_{41}]$&4\\
$[T_{41},1/2]$&1\\
\bottomrule
\end{tabular}
\end{center}
The relevant parabola centers are $-p,-u,0,t,2t,q$, with biases
$\eta,3\eta,0,2\eta,4\eta,\eta$. Subtracting neighboring
quadratics gives exactly \eqref{eq:phase-boundaries}. The remaining
copies are above this envelope: their differences are affine, so it
suffices to check the displayed interval endpoints. The successive
active widths are $\omega,p-\omega,t,\omega$, with the two outer
pieces totaling $p-\omega$. Positivity follows, in particular, from
$7773^2-10(2447)^2=541439>0$.

For completeness all $j\ge5$ have a uniform positive price gap. Using
only $j=0,1,2$, the relevant centers are $-p,0,t,q$ and the intersections
are $-p/2-\eta/(2p)$, $t/2+\eta/t$ and
$(q+t)/2-\eta/(2p)$. The largest potentially troublesome intersection
price is at most $2\eta+p^2/4<5\eta$, with gap greater than
$59/10000$, using
\[
 4/25<t<1/6,\qquad5/12<p<21/50,\qquad1/60<\eta<1/55.
\]
The other intersection prices are below $1/16<5\eta$;
the endpoint price $\eta+t^2/4$ is also below $5\eta$.
But $f_j\ge5\eta$ for $j\ge5$. The bounded-$j$ comparison
therefore gives $j_N\le4$ eventually.

Away from $O(N^{-1})$ neighborhoods of the five crossings, the actual
winner has the displayed code. Since $p$ is irrational, its rotation
is equidistributed. Those shrinking neighborhoods have density zero:
their upper density is bounded by the length of any fixed larger
neighborhood. The interval widths prove the claimed padding frequencies.

\subsection{Boundary drift and exact eventual record gaps}

To decide every sufficiently late crossing, extend
Lemma~\ref{lem:phase-uniform} to the seventh CDF order. Its smooth
efficiency expansion has remainder $O(m^{-4})$, including the branch
derivatives needed for optimization. The seventh-order coefficient's
value at $(p,c_*)$ will be denoted $J_*$. Its value is common to the
core comparisons and will cancel.

All derivatives below are of the explicit $Q,H$ above at $(p,c_*)$.
Put
\begin{align*}
 K(e)&=(2\alpha+Q_{pc})e+Q_p+H_c,\qquad
 F_2(e)=\alpha e^2+Q_pe+H_0,\\
 G(e)&=E_3e^3-Q_{pp}e^2/2-H_pe+K(e)^2/(2Q_{cc}).
\end{align*}
The needed constants are
\[
\begin{gathered}
 Q_{cc}=\phz(35+11s)/15,\quad
 Q_{pc}=\phz(415+131s)/540,\quad
 Q_p=\phz(1483+463s)/19440,\\
 E_3=E'''(p)/6=-\phz(19+6s)/15,\quad
 c_0'(p)=-(23+4s)/108,\quad
 Q_{pp}=Q_{\min}''(p)+Q_{cc}c_0'(p)^2.
\end{gathered}
\]
The derivatives $H_c,H_p$ are obtained by differentiating the displayed
$T_5$ expression for $H$; thus every coefficient here is explicit.
Writing $c=c_*+d_1/m$ and $v=p+e/m+d_1/m^2$, branch stationarity gives
$d_1=-K(e)/Q_{cc}$. Substituting into the branch value yields
\begin{equation}\label{eq:phase-cubic}
 B_m=\ce-Q_*/m-F_2(e)/m^2+(G(e)-J_*)/m^3+O(m^{-4})
\end{equation}
on a winning integer branch. Near a core's internal half-integer
crossing both branches must be compared. At each policy crossing in
\eqref{eq:phase-boundaries}, the two competing cores' internal phases
lie in $[T_{30},T_{02}]\subset(-1/2,1/2)$, so there is a fixed distance
from those internal crossings.

For $\theta_j=\langle\theta+jp\rangle$ and $e_j=b+\theta_j$, put
\[
 L_j(\theta)=Q_*j+F_2(e_j),\qquad
 M_j(\theta)=-Q_*j^2-2jF_2(e_j)+G(e_j).
\]
Expanding \eqref{eq:phase-cubic} at $m=N-j$ gives
\[
 B_{N-j}=\ce-Q_*/N-L_j(\theta)/N^2
                    +(M_j(\theta)-J_*)/N^3+O(N^{-4}).
\]
At a transverse limiting crossing $T$ between $i,j$, the price
difference changes sign at
\[
 \theta=T+\frac{D_{ij}}N+O(N^{-2}),\qquad
 D_{ij}=\frac{M_i(T)-M_j(T)}{L_i'(T)-L_j'(T)}.
\]
Here the equation refers to the uniform sign comparison of branch
values; no continuous interpolation of integer-threshold optima is required.
Substitution of the explicit polynomials gives
\begin{equation}\label{eq:phase-drift}
\begin{gathered}
 D_{13}=D_{02}=D_{24}=D_A
   =\frac{129918319}{529079040}
       +\frac{191267395}{1693052928}s,\\
 D_{30}=D_{41}=D_B
   =-\frac{663979357}{2351462400}
       -\frac{20072514361}{141087744000}s.
\end{gathered}
\end{equation}
Both rational parts are nonzero, so neither $sD_A$ nor $sD_B$ is rational.

The actual rotation cannot remain ambiguously close to a shifted
boundary. Let $h_N=\theta_N-T$ and $\gamma=\delta-T$.
With $L_0=25920$, the numerators of
$\gamma=(\gamma_A+\gamma_Bs)/L_0$ are
\begin{center}
\begin{tabular}{crr}
\toprule
$T$&$\gamma_A$&$\gamma_B$\\
\midrule
$T_{13}$&95720&-25573\\
$T_{30}$&72401&-18232\\
$T_{02}$&43880&-12613\\
$T_{24}$&121640&-38533\\
$T_{41}$&98321&-31192\\
\bottomrule
\end{tabular}
\end{center}
These follow directly from
$\eta/t=(60+269s)/8640$ and $3\eta/(2u)=(269+6s)/2880$.
If $k$ is the nearest integer to $Np-\delta$, then
\[
 h_N=(A_N+B_Ns)/L_0,\quad
 A_N=L_0(k-2N)+\gamma_A,\quad B_N=12960N+\gamma_B.
\]
Its field norm is
\[
 \nu_N=(A_N^2-10B_N^2)/L_0^2\in L_0^{-2}\Z.
\]
In any fixed corridor $|h_N|\le C/N$, the conjugate of $h_N$ is
$-sN+O(1)$, so $\nu_N$ is bounded and takes only finitely many rational
values. Moreover $Nh_N=-\nu_N/s+O(N^{-1})$.
For the corresponding $D_A$ or $D_B$, each $-\nu/s-D$ is nonzero by
\eqref{eq:phase-drift}. The finite set has a positive minimum separation.
Thus $|\theta_N-T-D/N|\ge\varepsilon/N$ eventually for some
$\varepsilon>0$. Outside a sufficiently wide $C/N$ corridor the same
separation is immediate. This margin dominates the $O(N^{-2})$
boundary uncertainty. Other policies have fixed limiting price gaps
near each crossing; away from crossings the unique leading winner
persists. Hence the exact sufficiently late code is the six-interval
table, shifting the three $A$ endpoints by $D_A/N$ and the two $B$
endpoints by $D_B/N$. In particular the winning core is unique.

Finally the endpoint translations
\[
 T_{02}-p=T_{13},\qquad T_{13}+q=T_{24},\qquad
 T_{30}+q=T_{41}
\]
preserve the drift group. Changing $N$ to $N+1$ changes $D/N$ only
by $O(N^{-2})$, below the norm separation. Interval transport under
$\theta_{N+1}=\langle\theta_N-p\rangle$ therefore gives exactly
\[
 0\longrightarrow1,\quad1\longrightarrow0\text{ or }2,\quad
 2\longrightarrow0\text{ or }3,\quad3\longrightarrow4,\quad
 4\longrightarrow0.
\]
The return lengths to a record are $2,3,5$. Their frequencies per
horizon are $u-\omega,t-\omega,\omega$, respectively, all positive;
this follows by translating the displayed intervals, with the vanishing
shifts leaving frequencies unchanged. Their sum is $p-\omega$, and
their gap-weighted sum is one. Each gap thus occurs infinitely often,
proving Theorem~\ref{thm:iid-padding}.

\section{Local replacement of nearby support diameters}\label{sec:spread}

In this section we prove Theorem~\ref{thm:spread}. Dividing all summands by
the variance-weighted mean diameter $h_0$, we set
\[
 q_i=1-p_i,\qquad r_i=p_i^2+q_i^2,\qquad
 R_0=\sum_{i=1}^n a_i r_i,\qquad\mu=\sum_{i=1}^n p_i,
 \qquad\ell=\sum_{i=1}^n a_i|e_i|\le\sqrt A\,\sigma,
\]
where $\sum_{i=1}^n a_i e_i=0$ and $A/2\le R_0\le A$. Writing
$K=\sum_{i=1}^n B_i$ and $T=\sum_{i=1}^n e_i(B_i-p_i)$, the centered sum is
$S=K-\mu+T$ with variance $A+\sigma^2$. Set $c_k=\Pp(K=k)$ and
$N_j(k)=\E[T^j\ind_{\{K=k\}}]$.

\subsection{Probability-free Poisson-binomial and weighted-moment bounds}

Uniformly over all probabilities $p_i\in(0,1)$ and all $k\in\Z$, as
$A\to\infty$,
\begin{align}
 c_k&=\frac{\phi(x_k)}{\sqrt A}+O(A^{-1}),\notag\\
 F_K(k)-\Phi(x_k)
   &=\frac{\phi(x_k)}{\sqrt A}
      \left\{\frac12+\frac{\gamma_3}{6}(1-x_k^2)\right\}
           +O(A^{-1}),                              \label{eq:pb-grid}\\
 x_k&=\frac{k-\mu}{\sqrt A},\qquad
 \gamma_3=\frac1A\sum_{i=1}^n a_i(1-2p_i)\in[-1,1],\notag
\end{align}
and the unconditioned weighted moments $N_j(k)=\E[T^j\ind_{\{K=k\}}]$ satisfy
\begin{align}
 |N_1(k)|&\le C\ell A^{-3/2},\notag\\
 |N_2(k)-\sigma^2c_k|&\le \frac{C\sigma^2}{A},\notag\\
 0\le N_4(k)&\le \frac{C(\sigma^4+\eta^2\sigma^2)}{\sqrt A}.
                                                    \label{eq:weighted-coefficients}
\end{align}
To prove \eqref{eq:weighted-coefficients} without dividing by individual
characteristic functions, write $f_i(t)=q_i+p_i e^{\ii t}$, $s=e^{\ii t}-1$,
$\Psi=\prod_{i=1}^nf_i$, and $\Psi_i,\Psi_{ij}$ for the products with one or
two factors omitted. Because $\sum_{i=1}^n a_ie_i=0$, the Fourier transform
of $N_1$ is $\widehat N_1(t)=s\sum_i a_ie_i\Psi_i=-s^2\sum_i a_ip_ie_i\Psi_i$,
whose $L^1([-\pi,\pi])$ norm is $O(\ell A^{-3/2})$. Similarly, using
$\Psi_{ij}=\Psi+(1-f_i)\Psi_i+(1-f_j)\Psi_j+(1-f_i)(1-f_j)\Psi_{ij}$ and
$u_i=a_ie_i$ ($\sum_i u_i=0$) bounds $\|\widehat N_2-\sigma^2\Psi\|_1=O(\sigma^2/A)$,
and expanding $T^4$ into index partitions $(4),(3,1),(2,2),(2,1,1),(1,1,1,1)$
bounds $\|\widehat N_4\|_1=O((\sigma^4+\eta^2\sigma^2)/\sqrt A)$.

\subsection{Threshold localization prior to conditioning}

At any threshold $y=k-\mu+\theta$ with $k=\operatorname{round}(y+\mu)$ and
$|\theta|\le1/2$, decomposing $\{S\le y\}$ by the value of $K$ and summing
$N_4(j)/(|j-k|-1/2)^4$ over $j\ne k$ gives
\begin{equation}\label{eq:crossing-decomposition}
 F_h(y)=F_K(k)-c_k\Pp(T>\theta\mid K=k)+E_{k,\theta},
 \qquad |E_{k,\theta}|\le
             \frac{C(\sigma^4+\eta^2\sigma^2)}{\sqrt A},
\end{equation}
while the third-moment sum $R_h$ satisfies
\begin{equation}\label{eq:resource-change}
 R_h-R_0=3\sum_{i=1}^n a_i r_i e_i+3\sum_{i=1}^n a_i r_i e_i^2
                                     +\sum_{i=1}^n a_i r_i e_i^3,
 \qquad |R_h-R_0|\le C(\ell+\sigma^2+\eta\sigma^2).
\end{equation}
Thus any threshold achieving $\mathcal C_h(y)\ge c_0$ has bounded $|x_k|$,
so $c_k\ge c_{c_0}/\sqrt A$, and dividing \eqref{eq:weighted-coefficients}
by $c_k$ yields
\begin{align}
 |\E(T\mid K=k)|&\le \frac{C\ell}{A},\notag\\
 \E(T^2\mid K=k)&=\sigma^2\bigl\{1+O(A^{-1/2})\bigr\},\notag\\
 \E(T^4\mid K=k)&\le C(\sigma^4+\eta^2\sigma^2).
                                                    \label{eq:conditional-bulk}
\end{align}

\subsection{Conditional $L^1$ lower bound via rejective-sampling matchings}

We next prove that at any localized count $k$,
\begin{equation}\label{eq:conditional-l1}
                   \E(|T|\mid K=k)\ge \frac{c\ell}{\sqrt A}.
\end{equation}
The conditional law of $(B_1,\ldots,B_n)$ given $K=k$ is a rejective
sampling law \citep[Section~2]{BoistardLopuhaaRuizGazen2012}; by
\eqref{eq:pb-grid}, for any set $F$ of at most three indices,
$\Pp(B_F=b_F\mid K=k)=\Pp(B_F=b_F)\{1+O(A^{-1/2})\}$. Let $B,B'$ be two
independent draws from the conditional law given $K=k$, and let
$C=\{i:B_i\ne B_i'\}$, which has even cardinality $m=|C|$. Conditional on
$C$ and on the common values outside $C$, the signs $B_i-B_i'\in\{-1,1\}$ for
$i\in C$ are uniformly distributed over balanced sign vectors (since the
odds product $\prod_{i\in C}(p_i/q_i)$ is independent of which half of $C$ is
$+1$). Generating a uniform balanced sign vector by choosing a uniform
perfect matching of $C$ and assigning independent Rademacher signs to each
matched pair gives, by Khintchine's inequality on each matching,
\[
 \E(|T-T'|\mid C)
 \ge\sqrt{\frac23}\,
       \frac{\sum_{i<j\in C}|e_i-e_j|}{(m-1)\sqrt m}\qquad(m\ge2).
\]
Since $\Pp(i,j\in C)\ge ca_ia_j$ and $\E(m\mid i,j\in C)\le CA$, Jensen's
inequality for $m^{-3/2}$ and $\sum_ja_j|e_i-e_j|\ge A|e_i|$ yield
\[
 2\E(|T|\mid K=k)\ge\E(|T-T'|\mid K=k)
 \ge cA^{-3/2}\sum_{i<j}a_ia_j|e_i-e_j|
 \ge \frac{c\ell}{\sqrt A},
\]
proving \eqref{eq:conditional-l1}.

\subsection{Completion of the comparison and optimality of the rate}

Setting $M_1=\E(|T|\mid K=k)$, $W=\Pp(T>\theta\mid K=k)$, and $L=\theta+W$,
equations \eqref{eq:conditional-bulk}--\eqref{eq:conditional-l1} and Cauchy--Schwarz
give $|\E(T\mid K=k)|\le M_1/4$, $M_1\ge c\sigma^2/(\sigma+\eta)$, and
\begin{equation}\label{eq:relative-phase}
 L\ge c\max\left(\frac\ell{\sqrt A},\,
                      \frac{\sigma^2}{\sigma+\eta}\right)>0,
 \qquad |\theta|+\theta^2+W\le CL,
 \qquad\sigma^2\le C(\sigma+\eta)L.
\end{equation}
Comparing $D_h(y)=F_h(y)-\Phi(y/\sqrt{A+\sigma^2})$ directly with the
unperturbed discrepancy $D_0(k)=F_K(k)-\Phi(x_k)$ via \eqref{eq:crossing-decomposition}
and \eqref{eq:resource-change} shows that every error term is $O((\sigma+\eta)^3L+L/\sqrt A)$,
giving
\[
 \mathcal C_h(y)-\frac{A^{3/2}}{R_0}D_0(k)
 \le\bigl\{-c+C(\sigma+\eta)^3+CA^{-1/2}\bigr\}L
 \le-\kappa\frac{\sigma^2}{\sigma+\eta},
\]
which proves Theorem~\ref{thm:spread} (and when $p_i\in J\subset(0,1)$ for a
compact set $J$, $\eta\le\sigma/\sqrt{\min_{p\in J}p(1-p)}$ gives
$\sigma^2/(\sigma+\eta)\ge c_J\sigma$).

Finally, to see that the rate $\sigma^2/(\sigma+\eta)$ cannot be improved to
$\sigma$ without a lower bound on $p_i$, take a fixed large i.i.d.\ Esseen
group of size $m$ with ratio $>c_0$ attained on the positive closed branch,
and adjoin one rare coordinate $B_{m+1}\sim\Bern(s)$ ($s\to0$) with diameter
$g_s=1-\eta$ ($0<\eta<1/4$) while giving the $m$ group coordinates diameter
$g_c=1+s(1-s)\eta/(ma)$. Because $0<g_s/g_c<1$, the closed CDF at uncentered
threshold $kg_c$ equals $(1-s)F_{\mathrm{group}}(k)+sF_{\mathrm{group}}(k-1)$,
identical to the common-diameter value at $\eta=0$. Since the smooth branch
ratio is analytic in $(s,\eta)$ near $(0,0)$ and is independent of $\eta$ at
$s=0$, its mixed partial derivative is bounded, so
$0<\mathcal C_{\mathrm{com}}(s)-\mathcal C_h\le K_m s\eta$. Choosing $s=\eta^4$
gives $\sigma\sim\sqrt s\,\eta=\eta^3$ and gain $O(\eta^5)=O(\sigma^2/(\sigma+\eta))$.

\bibliographystyle{plainnat}
\bibliography{refs}
\end{document}
